\chapter*{Resumen}
\addcontentsline{toc}{chapter}{Resumen}

ImprontIA es una plataforma de apoyo a la integridad académica en evaluaciones digitales. Aborda la escasa visibilidad sobre el proceso con el que se produce una respuesta en una evaluación remota, frente a una supervisión audiovisual percibida como invasiva. Aporta trazabilidad de ese proceso mediante señales conductuales de interacción ---dinámica de tipeo, movimiento del cursor, cambios de foco y pegado---, sin cámara ni micrófono y sin incorporar el contenido escrito a la telemetría conductual ni al modelo de anomalías (la respuesta final se conserva por separado como entrega académica): mide la atipicidad de la interacción respecto de la cohorte de la misma evaluación y genera indicadores trazables que un docente revisa; la decisión es siempre humana.

El trabajo se fundamenta en un marco teórico de nueve ejes y en un estado del arte sobre \textit{keystroke dynamics}, \textit{mouse dynamics}, \textit{proctoring} y detección de asistencia algorítmica. Una investigación con usuarios (encuesta a 213 estudiantes y tres entrevistas) fijó las condiciones de aceptabilidad: transparencia, revisión humana, proporcionalidad y separación entre alerta y decisión académica. El producto mínimo viable (MVP), desarrollado de forma iterativa e incremental, consta de un \textit{frontend} React con captura embebida y consentimiento previo, y de un \textit{backend} FastAPI que genera indicadores por reglas al finalizar cada sesión y aplica \textit{Isolation Forest} sobre la cohorte completa al cerrar la evaluación, con 178 pruebas automatizadas aprobadas.

La validación técnica combinó un conjunto de datos sintético con etiquetas conocidas, generado por el equipo, y pruebas automatizadas sobre el sistema real. Con el protocolo vigente de cohorte al cierre, el modelo aportó detecciones adicionales exclusivamente en el perfil sintético de ritmo distinto, con una detección del 54 al 67\,\% según la composición y el tamaño de la cohorte; estos resultados verifican el comportamiento del sistema en condiciones sintéticas y no estiman su desempeño con estudiantes. Las pruebas de punta a punta dieron el resultado esperado en 47 de 48 sesiones, el sistema cumplió las cotas de tiempo de procesamiento y de consulta, excepto con 200 entregas simultáneas, y una evaluación experta de usabilidad llevó a cero las tareas con problemas y las fallas serias de accesibilidad.

La viabilidad legal, técnica y económica es factible: el tratamiento se ajusta a la ley de protección de datos personales con consentimiento, minimización y plazos de conservación acotados, la infraestructura no requiere \textit{hardware} especial y la suscripción anual por institución es viable en el escenario base, con una tasa interna de retorno del 23,8\,\% frente al 20,36\,\% exigido. El siguiente paso es un piloto con voluntarios que calibre el sistema con datos reales.

\vspace{4pt}
\noindent\textbf{Palabras clave:} trazabilidad de autoría; biometría conductual; \textit{keystroke dynamics}; integridad académica; detección de anomalías; \textit{Isolation Forest}; privacidad por diseño; revisión humana.

\newpage
\chapter*{Abstract}
\addcontentsline{toc}{chapter}{Abstract}
\begin{otherlanguage}{english}

ImprontIA is a platform that supports academic integrity in digital assessments. It addresses the limited visibility of the process through which an answer is produced in a remote assessment, compared with audiovisual supervision perceived as invasive. It provides traceability of that process through behavioral interaction signals ---typing dynamics, cursor movement, focus changes and pasting---, without a camera or microphone and without incorporating the written content into the behavioral telemetry or the anomaly model (the final answer is kept separately as the academic submission): it measures the atypicality of the interaction relative to the cohort of the same assessment and produces traceable indicators that a teacher reviews; the decision is always human.

The work is grounded on a theoretical framework of nine axes and on a state of the art covering keystroke dynamics, mouse dynamics, proctoring and detection of algorithmic assistance. User research (a survey of 213 students and three interviews) established the acceptability conditions: transparency, human review, proportionality and separation between alert and academic decision. The minimum viable product (MVP), developed iteratively and incrementally, consists of a React frontend with embedded capture and prior consent, and of a FastAPI backend that generates rule-based indicators when each session ends and applies Isolation Forest over the complete cohort when the assessment is closed, with 178 passing automated tests.

The technical validation combined a synthetic dataset with known labels, generated by the team, and automated tests on the real system. With the current cohort-at-close protocol, the model contributed additional detections only in the synthetic different-rhythm profile, detecting between 54 and 67\,\% depending on the composition and size of the cohort; these results verify the behavior of the system under synthetic conditions and do not estimate its performance with students. End-to-end tests produced the expected result in 47 of 48 sessions, the system met the processing and query time bounds, except with 200 simultaneous submissions, and an expert usability evaluation brought the tasks with problems and the serious accessibility failures down to zero.

Legal, technical and economic viability is feasible: the processing complies with Law~25,326 through consent, data minimization and limited retention periods, the infrastructure requires no special hardware and the annual subscription per institution is viable in the base scenario, with an internal rate of return of 23.8\,\% against the required 20.36\,\%. The next step is a pilot with volunteers to calibrate the system with real data.

\vspace{4pt}
\noindent\textbf{Keywords:} authorship traceability; behavioral biometrics; keystroke dynamics; academic integrity; anomaly detection; Isolation Forest; privacy by design; human review.
\end{otherlanguage}
